\documentclass[a4paper, 12pt]{article}
 \usepackage[utf8]{inputenc}
\usepackage[T1]{fontenc}
\usepackage[margin=1in]{geometry} 
\usepackage{amsmath,amsthm,amssymb,amsfonts, enumitem, fancyhdr, color, comment, graphicx, environ}
\usepackage{enumitem}




% configuración del fondo-------------------------------
\usepackage{eso-pic}
\newcommand\BackgroundPic{
\put(0,0){
\parbox[b][\paperheight]{\paperwidth}{%
\vfill
\centering
\includegraphics[width=\paperwidth,height=\paperheight,
keepaspectratio]{Logo.jpg}% 
\vfill
}}}





\usepackage{lastpage}
\usepackage{fancyhdr}

% PIE DE PAGINA ------------------------------------------
\fancyhf{} % Numeración en pie de pagina
\fancyfoot[c]{\thepage \hspace{1pt} de \pageref{LastPage}}
\addtolength{\textheight}{-0.5in}

% ENCABEZADO ---------------------------------------------
\pagestyle{fancy}
\setlength{\headheight}{65pt} %Separacion de la linea y el Titulo 
\renewcommand{\headrulewidth}{1pt}

% Texto a la Izquierda del encabezado -----
\lhead{
    Facultad de Ingeniería Económica, Estadística y Ciencias Sociales \\
    Escuela de Ingeniería Estadística\\
    \textbf{Ciclo:} 2024-I
} 

% Texto a la Derecha del encabezado -----
\rhead{     
    27/04/2024\\
    \textbf{Curso:} Nombre del Curso\\
    \textbf{Profesor:} Yhon Tiahuallpa
}  


% Puntaje en el margen izquierdo----------------------------
\usepackage{marginnote}
\newcommand{\pts}[1]{\reversemarginpar\marginnote{ (#1 pts) }}




% ________________________________________________________


\begin{document}

\AddToShipoutPicture{\BackgroundPic} % Agrega el fondo a todas las hojas

$$\text{\textbf{\large Examen Parcial/Final}}$$

\vspace{0.3cm}



\begin{enumerate}[labelwidth=0pt, align=left, leftmargin=*]

    \item \pts{2}  Sea $\left\{X_n\right\}$ un proceso estacionario en covarianza con media cero, con función de covarianza $R_X(v)$ y función de densidad espectral $f_X(\omega),-\pi \leq \omega \leq \pi$. 
    $$
    Y_n=\sum_{k=0}^{\infty} a_k X_{n-k} \;\;.
    $$
    Demuestre que la ... y....
        $$
    \begin{aligned}
        f_Y(\omega) & =\frac{\sigma_X^2}{\sigma_Y^2}\left|\sum_{k=0}^{\infty} a_k e^{i k \omega}\right|^2 f_X(\omega) \\
         &=\frac{\sigma_X^2}{\sigma_Y^2}\left|\sum_{k=0}^{\infty} a_k e^{i k \omega}\right|^2 f_X(\omega) \quad-\pi \leq \omega \leq \pi.
    \end{aligned}
    $$



    \item \pts{3} El puntaje puede ir antes del enunciado



    
    \item   Calcular la función de densidad espectral del proceso autorregresivo $\left\{X_n\right\}$ el cual  
    satisface
    
    $$
        (a+b)^2=a^2+2ab+b^2
    $$
        
    \begin{itemize}
        \item \pts{2} El puntaje puede ir por items

        $$f(\omega)=\frac{\sigma_X^2}{\sigma_Y^2}\left|\sum_{k=0}^{\infty} a_k e^{i k \omega}\right|^2 f_X(\omega)$$

        \item \pts{3} donde $\left\{\xi_n\right\}$ y ademas...  son todas menores que uno en valor absoluto.

        $$f(\omega)=\frac{\sigma_X^2}{\sigma_Y^2}\left|\sum_{k=0}^{\infty} a_k e^{i k \omega}\right|^2 f_X(\omega)$$
    \end{itemize} 
    
    \qquad \textit{\textbf{Solución:}}
    $$
        f(\omega)=\frac{\sigma_X^2}{\sigma_Y^2}\left|\sum_{k=0}^{\infty} a_k e^{i k \omega}\right|^2 f_X(\omega)
    $$
     
    
    
    
    
    
    \item \pts{5} Sea $\left\{X_n\right\}$ y $\alpha_0, \ldots, \alpha_q$ son reales y $\left\{\xi_n\right\}$, ademas de:

    $$
        f(\omega)=\frac{\sigma_X^2}{\sigma_Y^2}\left|\sum_{k=0}^{\infty} a_k e^{i k \omega}\right|^2 f_X(\omega)
    $$
     


    
    donde $z_1, \ldots, z_q$ son las $q$ raíces de
    
    $$
        f(\omega)=\frac{\sigma_X^2}{\sigma_Y^2}\left|\sum_{k=0}^{\infty} a_k e^{i k \omega}\right|^2 f_X(\omega)
    $$
     
    
    
    
    
    \item \pts{5} Sean $\left\{\xi_n\right\}$  variables aleatorias independientes idénticamente distribuidas con medias nulas y varianzas unitarias. Demostrar que todo proceso de media movil
        $$
        f(\lambda)=\frac{1}{2 \pi \sigma_X^2} \prod_{j=1}^q\left|e^{i \lambda}-z_j\right|^2,
        $$
    es ergódico. Supongamos que $\displaystyle \sum a_k^2<\infty$. ¿Es lo mismo para $\displaystyle Y_n=\sum_{k=0}^{\infty} a_k \xi_{n-k}$?.

        $$
        f(\lambda)=\frac{1}{2 \pi \sigma_X^2} \prod_{j=1}^q\left|e^{i \lambda}-z_j\right|^2,
        $$


        
    








\end{enumerate}






%%%%%%%%%%%%%%%%%%%%%%%%%%%%%%%%%%%%%%%%
\end{document}